% GENERATED FILE. Edit canonical lessons, then run npm run generate:books.
% canonical-source-hash: fnv1a64:69239e9747759a49
% canonical-lessons: ZH-C318-biyou, ZH-C318-xinfeng, ZH-C318-mingxinpian, ZH-C318-huiyin, ZH-C318-zhuangao

\chapter{Letters from Far Away}
\label{ch:zh-letters-from-far-away}
\hlchaptermodality{318}

\begin{chapteropening}
\textbf{By the end of this chapter:} \emph{I can talk about a pen pal, an envelope and a postcard, wait for a reply and pass a message on.}

Complete the last lesson of chapter 318: I can talk about a pen pal, an envelope and a postcard, wait for a reply and pass a message on.
\end{chapteropening}

\section[bǐyǒu]{bǐyǒu — pen pal}

\label{lesson:ZH-C318-biyou}



\begin{quote}
Before the new one: say the Mandarin for household, then the Mandarin for beforehand.
\end{quote}

\subsection*{You'll want to know: bǐyǒu}
\emph{bǐyǒu} — \textquotedblleft{}pen pal\textquotedblright{}.

\begin{quote}
\emph{wǒ xiǎoshíhou yǒu yí ge bǐyǒu} — When I was little I had a pen pal.
\end{quote}

Its characters come later; for now it is a word for the ear.

Two third tones meet in this word: bǐ is said with a rising tone before yǒu, and the spelling keeps both marks.

\begin{cousinweb}
\emph{bǐyǒu} is built from \emph{bǐ} 'a pen' and \emph{yǒu} 'a friend', both inherited from Middle Chinese, \emph{pit} and \emph{hjuwX} (Guangyun, in Baxter's transcription), which names the first as the writing brush and glosses the second as those who share one mind.
\end{cousinweb}

\begin{grammarlens}[title={What you've built}]
One more word for telling what happened.
\end{grammarlens}

\subsection*{Your turn}
\begin{itemize}
\raggedright
  \item \emph{Say it:} \emph{bǐyǒu}
  \item \emph{Say it:} \emph{bǐyǒu}, once more
  \item \emph{Say it:} \emph{wǒ xiǎoshíhou yǒu yí ge bǐyǒu} — When I was little I had a pen pal.
  \item \emph{Recall:} say the Mandarin for household, then the Mandarin for beforehand, then say \emph{bǐyǒu} again
\end{itemize}

\begin{tcolorbox}[breakable,colback=teal!4,colframe=teal!35!black,title={Before you move on}]
What does bǐyǒu mean? (\textquotedblleft{}Pen pal\textquotedblright{}.) Say it once more.
\end{tcolorbox}
\section[xìnfēng]{xìnfēng — envelope}

\label{lesson:ZH-C318-xinfeng}



\begin{quote}
Before the new one: say the Mandarin for beforehand, then the Mandarin for a pen pal.
\end{quote}

\subsection*{You'll want to know: xìnfēng}
\emph{xìnfēng} — \textquotedblleft{}envelope\textquotedblright{}.

\begin{quote}
\emph{tā de xìnfēng lǐ yǒu yì zhāng zhàopiàn} — There was a photo in his envelope.
\end{quote}

Its characters come later; for now it is a word for the ear.

\begin{cousinweb}
\emph{xìnfēng} is built from \emph{xìn} 'a letter' and \emph{fēng} 'to seal', both inherited from Middle Chinese, \emph{sinH} and \emph{pjowng} (Guangyun, in Baxter's transcription). The dictionary glosses the first 'loyal and true' and the second 'great', 'a domain' and 'a rank'; the senses 'a letter' and 'to seal' are not given there.
\end{cousinweb}

\begin{grammarlens}[title={What you've built}]
One more word for telling what happened.
\end{grammarlens}

\subsection*{Your turn}
\begin{itemize}
\raggedright
  \item \emph{Say it:} \emph{xìnfēng}
  \item \emph{Say it:} \emph{xìnfēng}, once more
  \item \emph{Say it:} \emph{tā de xìnfēng lǐ yǒu yì zhāng zhàopiàn} — There was a photo in his envelope.
  \item \emph{Recall:} say the Mandarin for beforehand, then the Mandarin for a pen pal, then say \emph{xìnfēng} again
\end{itemize}

\begin{tcolorbox}[breakable,colback=teal!4,colframe=teal!35!black,title={Before you move on}]
What does xìnfēng mean? (\textquotedblleft{}Envelope\textquotedblright{}.) Say it once more.
\end{tcolorbox}
\section[míngxìnpiàn]{míngxìnpiàn — postcard}

\label{lesson:ZH-C318-mingxinpian}



\begin{quote}
Before the new one: say the Mandarin for a pen pal, then the Mandarin for an envelope.
\end{quote}

\subsection*{You'll want to know: míngxìnpiàn}
\emph{míngxìnpiàn} — \textquotedblleft{}postcard\textquotedblright{}.

\begin{quote}
\emph{tā cóng Zhōngguó gěi wǒ jì le yì zhāng míngxìnpiàn} — He sent me a postcard from China.
\end{quote}

Its characters come later; for now it is a word for the ear.

\begin{cousinweb}
\emph{míngxìnpiàn} is built from \zh{明} 'bright, open', \emph{xìn} 'a letter' and \emph{piàn} 'a flat piece', all three inherited from Middle Chinese, \emph{mjaeng}, \emph{sinH} and \emph{phenH} (Guangyun, in Baxter's transcription), which glosses the first 'light' and the last 'half' and 'a split piece of wood'.
\end{cousinweb}

\begin{grammarlens}[title={What you've built}]
One more word for telling what happened.
\end{grammarlens}

\subsection*{Your turn}
\begin{itemize}
\raggedright
  \item \emph{Say it:} \emph{míngxìnpiàn}
  \item \emph{Say it:} \emph{míngxìnpiàn}, once more
  \item \emph{Say it:} \emph{tā cóng Zhōngguó gěi wǒ jì le yì zhāng míngxìnpiàn} — He sent me a postcard from China.
  \item \emph{Recall:} say the Mandarin for a pen pal, then the Mandarin for an envelope, then say \emph{míngxìnpiàn} again
\end{itemize}

\begin{tcolorbox}[breakable,colback=teal!4,colframe=teal!35!black,title={Before you move on}]
What does míngxìnpiàn mean? (\textquotedblleft{}Postcard\textquotedblright{}.) Say it once more.
\end{tcolorbox}
\section[\texorpdfstring{\zh{回音} (huíyīn)}{回音 (huíyīn)}]{\zh{回音} (huíyīn) — reply; echo}

\label{lesson:ZH-C318-huiyin}



\begin{quote}
Before the new one: say the Mandarin for an envelope, then the Mandarin for a postcard.
\end{quote}

\subsection*{You'll want to know: \zh{回音}}
\textbf{\zh{回音}} — \emph{huíyīn} — \textquotedblleft{}reply; echo\textquotedblright{}.

\begin{quote}
\emph{wǒ gěi tā xiě le xìn tā réngrán méiyǒu huíyīn} — I wrote to him, and there's still no reply.
\end{quote}

\begin{cousinweb}
\textbf{\zh{回音}} (\emph{huíyīn}) is built from \zh{回} 'to turn back' and \zh{音} 'a sound', both inherited from Middle Chinese, \emph{hwoj} and \emph{'im} (Guangyun, in Baxter's transcription), which glosses the first 'to turn' and the second 'a sound'. A sound that comes back is an echo; the sense 'a reply' is not given there.
\end{cousinweb}

\begin{grammarlens}[title={What you've built}]
One more word for telling what happened.
\end{grammarlens}

\subsection*{Your turn}
\begin{itemize}
\raggedright
  \item \emph{Say it:} \emph{huíyīn}
  \item \emph{Say it:} \emph{huíyīn}, once more
  \item \emph{Say it:} \emph{wǒ gěi tā xiě le xìn tā réngrán méiyǒu huíyīn} — I wrote to him, and there's still no reply.
  \item \emph{Recall:} say the Mandarin for an envelope, then the Mandarin for a postcard, then say \emph{huíyīn} again
\end{itemize}

\begin{tcolorbox}[breakable,colback=teal!4,colframe=teal!35!black,title={Before you move on}]
What does \zh{回音} mean? (\textquotedblleft{}Reply; echo\textquotedblright{}.) Say it once more.
\end{tcolorbox}
\section[zhuǎngào]{zhuǎngào — to pass on a message}

\label{lesson:ZH-C318-zhuangao}



\begin{quote}
Before the new one: say the Mandarin for a postcard, then the Mandarin for a reply to a letter.
\end{quote}

\subsection*{You'll want to know: zhuǎngào}
\emph{zhuǎngào} — \textquotedblleft{}to pass on a message\textquotedblright{}.

\begin{quote}
\emph{qǐng nǐ zhuǎngào tā míngtiān wǒ bú qù le} — Please tell him for me that I'm not going tomorrow.
\end{quote}

Its characters come later; for now it is a word for the ear.

\begin{cousinweb}
\emph{zhuǎngào} is built from \emph{zhuǎn} 'to turn' and \emph{gào} 'to tell', both inherited from Middle Chinese, \emph{trjwenX} and \emph{kawH} (Guangyun, in Baxter's transcription), which glosses the first 'to move' and the second 'to report'.
\end{cousinweb}

\begin{grammarlens}[title={What you've built}]
One more word for telling what happened.
\end{grammarlens}

\subsection*{Your turn}
\begin{itemize}
\raggedright
  \item \emph{Say it:} \emph{zhuǎngào}
  \item \emph{Say it:} \emph{zhuǎngào}, once more
  \item \emph{Say it:} \emph{qǐng nǐ zhuǎngào tā míngtiān wǒ bú qù le} — Please tell him for me that I'm not going tomorrow.
  \item \emph{Recall:} say the Mandarin for a postcard, then the Mandarin for a reply to a letter, then say \emph{zhuǎngào} again
\end{itemize}

\begin{tcolorbox}[breakable,colback=teal!4,colframe=teal!35!black,title={Before you move on}]
What does zhuǎngào mean? (\textquotedblleft{}To pass on a message\textquotedblright{}.) Say it once more.
\end{tcolorbox}
